% Auto-generated by render_prompts.py -- verbatim XLRS-Bench-lite items
\begin{table}[H]
\centering
\caption{Example XLRS-Bench questions (verbatim). R marks tasks in the relational group, N the non-relational group.}
\label{tab:examples}
\scriptsize
\setlength{\tabcolsep}{3pt}
\begin{tabular}{@{}p{0.22\linewidth}c p{0.70\linewidth}@{}}
\toprule
Task & Grp & Question (answer) \\
\midrule
Counting with Complex Reasoning & R & How many parking lots are there in the town shown in the picture above? \emph{(D: 9)} \\
Object Spatial Relationship & R & In the image, the two smaller green rectangular greenhouses located towards the upper right are positioned in which direction relative to the larger multi-unit residential area? \emph{(C: To the right)} \\
Regional Counting & R & How many houses are there inside the red circle? \emph{(B: 5)} \\
Route Planning & R & In the image below, what is the shortest driving route from the building with the red pyramid roof to the parking lot at the port on the left? \emph{(B: Exit the building and head west. Continue straight through the first intersection. At the next intersection, turn left to reach your destination.)} \\
Anomaly Detection and Interpretation & N & The width and calmness of the river in the picture suggest? \emph{(D: It may serve as a vital natural water source, crucial for the surrounding ecosystem.)} \\
Object Classification & N & Identify the object category within the given bounding box in the satellite or aerial image. Image resolution: 4096 x 4097. Bounding box: [2368, 2061, 2447, 2108]. \emph{(B: The Red Roof)} \\
Object Color & N & Determine the color of the object within the given reference bounding box in the satellite image. Image resolution: 5616 x 3744. Bounding box: [3826, 1648, 3944, 1695]. \emph{(C: White)} \\
Overall Land Use Classification & N & Select the options that best represent the types of land use depicted in the image. \emph{(ABD: Farmland; Lakes or Ponds; Parking Lot or Garage [multi-select])} \\
Regional Land Use Classification & N & Select the type of land marked with a red circle in the image. \emph{(D: Farmland)} \\
\bottomrule
\end{tabular}
\end{table}
